% GENERATED FILE. Edit canonical lessons, then run npm run generate:books.
% canonical-source-hash: fnv1a64:fe6ff38dcf6740d7
% canonical-lessons: RU-C52-book, RU-C52-pen, RU-C52-paper, RU-C52-notebook, RU-C52-pencil

\chapter{Book, Pen, Paper, Notebook, Pencil}
\label{ch:ru-book-pen-paper-notebook-pencil}
\hlchaptermodality{52}

\begin{chapteropening}
\textbf{By the end of this chapter:} \emph{I can name a book, a pen, paper, a notebook and a pencil.}

Complete the last lesson of chapter 52: I can name a book, a pen, paper, a notebook and a pencil.
\end{chapteropening}

\section[\texorpdfstring{\ru{книга} (kníga)}{книга (kníga)}]{\ru{книга} (kníga) — a book}

\label{lesson:RU-C52-book}



\begin{quote}
Before the new one: say the Russian for a knife, then the Russian for a fork.
\end{quote}

\subsection*{You'll want to know: \ru{книга}}
\textbf{\ru{книга}} (\emph{kníga}) — \textquotedblleft{}a book\textquotedblright{}.

\begin{cousinweb}
\textbf{\ru{книга}} $\leftarrow$ Proto-Slavic \emph{*kŭniga}. Vasmer, following Räsänen, takes it as a loan through early Turkic, Danube Bulgar \emph{küiniv}, ultimately from Chinese \emph{k'üen}, \textquotedblleft{}a scroll\textquotedblright{}; Hungarian \emph{könyv}, \textquotedblleft{}book\textquotedblright{}, has the same source. He finds the older Assyrian account less likely.
\end{cousinweb}

\begin{grammarlens}[title={What you've built}]
One more word for this chapter.
\end{grammarlens}

\subsection*{Your turn}
\begin{itemize}
\raggedright
  \item \emph{Say it:} \emph{kníga}
  \item \emph{Say it:} \emph{kníga}, once more
  \item \emph{Say it:} say \emph{kníga}
  \item \emph{Recall:} say the Russian for a knife, then the Russian for a fork, then say \emph{kníga} again
\end{itemize}

\begin{tcolorbox}[breakable,colback=teal!4,colframe=teal!35!black,title={Before you move on}]
What does \ru{книга} mean? (\textquotedblleft{}A book\textquotedblright{}.) Say it once more.
\end{tcolorbox}
\section[\texorpdfstring{\ru{ручка} (rúchka)}{ручка (rúchka)}]{\ru{ручка} (rúchka) — a pen}

\label{lesson:RU-C52-pen}



\begin{quote}
Before the new one: say the Russian for a fork, then the Russian for a book.
\end{quote}

\subsection*{You'll want to know: \ru{ручка}}
\textbf{\ru{ручка}} (\emph{rúchka}) — \textquotedblleft{}a pen\textquotedblright{}. Literally \textquotedblleft{}a little hand\textquotedblright{}.

\begin{cousinweb}
\textbf{\ru{ручка}} is \emph{\ru{рука}}, \textquotedblleft{}hand\textquotedblright{}, with \emph{-\ru{ка}}, \textquotedblleft{}little\textquotedblright{}; \emph{\ru{к}} turns to \emph{\ru{ч}} before it. The word means \textquotedblleft{}a little hand\textquotedblright{}, then \textquotedblleft{}a handle\textquotedblright{}, then \textquotedblleft{}a pen\textquotedblright{}, the thing you hold.
\end{cousinweb}

\begin{grammarlens}[title={What you've built}]
One more word for this chapter.
\end{grammarlens}

\subsection*{Your turn}
\begin{itemize}
\raggedright
  \item \emph{Say it:} \emph{rúchka}
  \item \emph{Say it:} \emph{rúchka}, once more
  \item \emph{Say it:} say \emph{rúchka}
  \item \emph{Recall:} say the Russian for a fork, then the Russian for a book, then say \emph{rúchka} again
\end{itemize}

\begin{tcolorbox}[breakable,colback=teal!4,colframe=teal!35!black,title={Before you move on}]
What does \ru{ручка} mean? (\textquotedblleft{}A pen\textquotedblright{}.) Say it once more.
\end{tcolorbox}
\section[\texorpdfstring{\ru{бумага} (bumága)}{бумага (bumága)}]{\ru{бумага} (bumága) — paper}

\label{lesson:RU-C52-paper}



\begin{quote}
Before the new one: say the Russian for a book, then the Russian for a pen.
\end{quote}

\subsection*{You'll want to know: \ru{бумага}}
\textbf{\ru{бумага}} (\emph{bumága}) — \textquotedblleft{}paper\textquotedblright{}.

\begin{cousinweb}
\textbf{\ru{бумага}} — Vasmer thinks it is most probably from Italian: Korsch derived it from \emph{bombagio}, \textquotedblleft{}cotton\textquotedblright{}, and most early paper in Russia was Italian. Greek \emph{bómbyx}, \textquotedblleft{}silk\textquotedblright{}, cannot explain the Russian form, he says.
\end{cousinweb}

\begin{grammarlens}[title={What you've built}]
One more word for this chapter.
\end{grammarlens}

\subsection*{Your turn}
\begin{itemize}
\raggedright
  \item \emph{Say it:} \emph{bumága}
  \item \emph{Say it:} \emph{bumága}, once more
  \item \emph{Say it:} say \emph{bumága}
  \item \emph{Recall:} say the Russian for a book, then the Russian for a pen, then say \emph{bumága} again
\end{itemize}

\begin{tcolorbox}[breakable,colback=teal!4,colframe=teal!35!black,title={Before you move on}]
What does \ru{бумага} mean? (\textquotedblleft{}Paper\textquotedblright{}.) Say it once more.
\end{tcolorbox}
\section[\texorpdfstring{\ru{тетрадь} (tetrád)}{тетрадь (tetrád)}]{\ru{тетрадь} (tetrád) — a notebook}

\label{lesson:RU-C52-notebook}



\begin{quote}
Before the new one: say the Russian for a pen, then the Russian for paper.
\end{quote}

\subsection*{You'll want to know: \ru{тетрадь}}
\textbf{\ru{тетрадь}} (\emph{tetrád}) — \textquotedblleft{}a notebook\textquotedblright{}.

\begin{cousinweb}
\textbf{\ru{тетрадь}} $\leftarrow$ Byzantine Greek \emph{tetrádion}, from \emph{tetrás}, \textquotedblleft{}a group of four\textquotedblright{}, a sheet folded into four leaves (Vasmer).
\end{cousinweb}

\begin{grammarlens}[title={What you've built}]
One more word for this chapter.
\end{grammarlens}

\subsection*{Your turn}
\begin{itemize}
\raggedright
  \item \emph{Say it:} \emph{tetrád}
  \item \emph{Say it:} \emph{tetrád}, once more
  \item \emph{Say it:} say \emph{tetrád}
  \item \emph{Recall:} say the Russian for a pen, then the Russian for paper, then say \emph{tetrád} again
\end{itemize}

\begin{tcolorbox}[breakable,colback=teal!4,colframe=teal!35!black,title={Before you move on}]
What does \ru{тетрадь} mean? (\textquotedblleft{}A notebook\textquotedblright{}.) Say it once more.
\end{tcolorbox}
\section[\texorpdfstring{\ru{карандаш} (karandásh)}{карандаш (karandásh)}]{\ru{карандаш} (karandásh) — a pencil}

\label{lesson:RU-C52-pencil}



\begin{quote}
Before the new one: say the Russian for paper, then the Russian for a notebook.
\end{quote}

\subsection*{You'll want to know: \ru{карандаш}}
\textbf{\ru{карандаш}} (\emph{karandásh}) — \textquotedblleft{}a pencil\textquotedblright{}.

\begin{cousinweb}
\textbf{\ru{карандаш}} $\leftarrow$ Turkic \emph{karadaš}, \textquotedblleft{}black stone\textquotedblright{}, as in Turkish \emph{karataş}, \textquotedblleft{}black slate\textquotedblright{} (Vasmer): \emph{kara}, \textquotedblleft{}black\textquotedblright{}, and \emph{taş}, \textquotedblleft{}stone\textquotedblright{}.
\end{cousinweb}

\begin{grammarlens}[title={What you've built}]
One more word for this chapter.
\end{grammarlens}

\subsection*{Your turn}
\begin{itemize}
\raggedright
  \item \emph{Say it:} \emph{karandásh}
  \item \emph{Say it:} \emph{karandásh}, once more
  \item \emph{Say it:} say \emph{karandásh}
  \item \emph{Recall:} say the Russian for paper, then the Russian for a notebook, then say \emph{karandásh} again
\end{itemize}

\begin{tcolorbox}[breakable,colback=teal!4,colframe=teal!35!black,title={Before you move on}]
What does \ru{карандаш} mean? (\textquotedblleft{}A pencil\textquotedblright{}.) Say it once more.
\end{tcolorbox}
